%% main.tex
%% "Coverage Maximization and Overlap Minimization for Multi-Ticket Lottery Selection"
%% Target: Electronic Journal of Combinatorics (EJC) — open access, free to publish
%%
%% Compile: pdflatex main && bibtex main && pdflatex main && pdflatex main

\documentclass[12pt]{article}

%% ── Encoding & fonts ──────────────────────────────────────────────────────────
\usepackage[T1]{fontenc}
\usepackage[utf8]{inputenc}
\usepackage{lmodern}
\usepackage{microtype}

%% ── Page geometry ─────────────────────────────────────────────────────────────
\usepackage[top=1.25in, bottom=1.25in, left=1.25in, right=1.25in]{geometry}
\setlength{\parskip}{4pt}
\linespread{1.08}

%% ── Mathematics ───────────────────────────────────────────────────────────────
\usepackage{amsmath, amsthm, amssymb, mathtools}

%% ── Theorem environments ──────────────────────────────────────────────────────
\theoremstyle{plain}
\newtheorem{theorem}{Theorem}[section]
\newtheorem{proposition}[theorem]{Proposition}
\newtheorem{lemma}[theorem]{Lemma}
\newtheorem{corollary}[theorem]{Corollary}
\newtheorem{conjecture}[theorem]{Conjecture}

\theoremstyle{definition}
\newtheorem{definition}[theorem]{Definition}
\newtheorem{example}[theorem]{Example}
\newtheorem{remark}[theorem]{Remark}

%% ── Algorithms ────────────────────────────────────────────────────────────────
\usepackage[ruled, linesnumbered, vlined]{algorithm2e}
\SetKwInput{KwInput}{Input}
\SetKwInput{KwOutput}{Output}

%% ── Tables & figures ──────────────────────────────────────────────────────────
\usepackage{bm}        % bold math symbols (\bm)
\usepackage{booktabs}
\usepackage{multirow}
\usepackage{array}
\usepackage{graphicx}
\usepackage{xcolor}
\usepackage{float}

%% ── Hyperlinks ────────────────────────────────────────────────────────────────
\usepackage[colorlinks=true, linkcolor=blue!60!black,
            citecolor=green!50!black, urlcolor=blue!70!black]{hyperref}

%% ── Bibliography ──────────────────────────────────────────────────────────────
\usepackage[numbers, sort&compress]{natbib}
\bibliographystyle{plainnat}

%% ── Custom macros ─────────────────────────────────────────────────────────────
% Universe, ticket size, prize threshold
\newcommand{\vv}{v}          % universe size
\newcommand{\kk}{k}          % ticket (block) size
% NOTE: \tt is reserved by LaTeX; use \tthresh for prize threshold
\newcommand{\tthresh}{t}     % prize threshold (hits required)
\newcommand{\nn}{n}          % number of tickets

% Key probability
\newcommand{\pone}{p(\vv,\kk,\tthresh)}   % single-ticket prize probability
\newcommand{\Pwin}{P_{\mathrm{win}}}  % P(at least one ticket wins)

% Covering number
\newcommand{\Cvkt}{C(\vv,\kk,\tthresh)}

% Hypergeometric distribution
\newcommand{\Hyp}[3]{\mathrm{Hyp}(#1,#2,#3)}

% Sets
\newcommand{\calA}{\mathcal{A}}   % ticket collection
\newcommand{\bbN}{\mathbb{N}}
\newcommand{\bbR}{\mathbb{R}}

% Shorthand for indicator
\newcommand{\ind}[1]{\mathbf{1}\!\left[#1\right]}

% Overlap
\newcommand{\overlap}[2]{|#1 \cap #2|}
% NOTE: \binom is already defined by amsmath — use it directly

%% ── Spacing tweaks ────────────────────────────────────────────────────────────
\renewcommand{\floatpagefraction}{0.8}
\renewcommand{\topfraction}{0.9}

%% ═══════════════════════════════════════════════════════════════════════════════
\begin{document}
%% ═══════════════════════════════════════════════════════════════════════════════

%% ── Title block ───────────────────────────────────────────────────────────────
\title{%
  Coverage Maximization and Overlap Minimization\\
  for Multi-Ticket Lottery Selection:\\
  Theory, Algorithms, and Empirical Validation
}

\author{%
  Lucas Rafael de Andrade\\[2pt]
  \small \textit{UERJ -- Programa de P\'{o}s-Gradua\c{c}\~{a}o em
    Ci\^{e}ncias Econ\^{o}micas (PPGCE)}\\
  \small \texttt{rafael.lucas@posgraduacao.uerj.br}
  \and
  Victor Hugo Nascimento\\[2pt]
  \small \textit{School of Applied Mathematics, FGV}\\
  \small Praia de Botafogo, Rio de Janeiro, Brazil\\
  \small \texttt{victorhugo.vhrn@gmail.com}
}

\date{\today}

\maketitle

%% ── Abstract ──────────────────────────────────────────────────────────────────
\begin{abstract}
We study the problem of selecting $n$ tickets in a $(v,k)$-lottery, where
each ticket is a $k$-subset of a universe $[v]$, a draw is a uniformly
random $k$-subset, and a prize is awarded for at least $t$ matches.  The
goal is to choose the collection of $n$ tickets that maximises the probability
that at least one ticket wins.

The classical combinatorial literature addresses a deterministic version of
this question, seeking ticket systems that guarantee a minimum number of
matches for every possible draw.  We take a probabilistic route instead and
show that, for two tickets, the probability of at least one win is a
non-increasing function of the pairwise overlap $|A_1 \cap A_2|$.  A
consequence, proved by exact enumeration for the Lotof\'{a}cil parameters
$(v=25, k=15, t=11)$, is that tickets with overlap at most six have mutually
exclusive winning events, so the joint prize probability reaches its maximum
value of $2p$ exactly.  For general $n \ge 2$ we conjecture the same
monotonicity, supported by simulation, and show that greedy coverage selection
followed by a simulated annealing refinement produces collections with
near-minimum pairwise overlap.  The greedy stage inherits the
$(1-1/e)$-approximation guarantee of \citet{Nemhauser1978}.

We apply both algorithms to the Lotof\'{a}cil, the first academic study of
this Brazilian lottery.  Using a paired Monte Carlo design with $50{,}000$
shared draws per condition, we estimate the joint prize probability for
$n = 1, 2, 3, 5, 7, 10, 15, 20, 25, 30, 40, 50, 75$, and~$100$ tickets.
The optimised strategy yields gains of up to 11 percentage points over
random selection at $n = 5$ (Cohen's $d = 2.64$, $p < 10^{-23}$), with
statistically significant advantages at every $n \ge 2$ up to the saturation
point near $n = 50$, beyond which both strategies achieve $P \ge 0.996$.
\end{abstract}

\vspace{4pt}
\noindent\emph{Keywords:} lottery design, covering design, set cover,
greedy algorithm, simulated annealing, hypergeometric distribution, pairwise
overlap, Lotof\'{a}cil.

\vspace{4pt}
\noindent\emph{MSC 2020:} 05B40 (Covering designs); 68W25 (Approximation
algorithms); 60C05 (Combinatorial probability).

\bigskip
\hrule
\bigskip

%% ── Sections (each in its own file) ──────────────────────────────────────────
%% sec1_intro.tex -- Section 1: Introduction

\section{Introduction}
\label{sec:intro}

A $(v,k)$-lottery is a game in which $k$ distinct numbers are drawn
uniformly at random from $[v] = \{1,\ldots,v\}$.  A player holding a ticket,
which is a $k$-subset of $[v]$, wins a prize when at least $t$ of the drawn
numbers appear on the ticket.  Games of this form are played in many countries
and have attracted sustained mathematical interest since the work of
Sch\"{o}nheim \cite{Schoenheim1964}, who introduced the study of covering
numbers $C(v,k,t)$, the minimum cardinality of a family of $k$-subsets
of $[v]$ such that every $t$-subset is contained in at least one.

The bulk of the literature pursues a deterministic goal: given $n$ tickets,
arrange them so that every possible draw $R$ satisfies $|A_i \cap R| \ge t$
for at least one $A_i$.  This is the $(v,k,t,m)$-lotto design framework
formalised in \cite{Furedi1996} and tabulated by Li and van Rees \cite{Li2002},
who describe greedy and simulated annealing constructions for small instances.
Gordon, Kuperberg and Patashnik \cite{Gordon1995} give efficient upper-bound
constructions; comprehensive tables appear in the \emph{Handbook of
Combinatorial Designs} \cite{Colbourn2007}.

The deterministic paradigm has a clear virtue: a well-chosen set of tickets
guarantees a prize regardless of the draw.  For practical budgets, however,
the number of tickets needed grows quickly with $v$ and $t$, and a player
committing to $n$ tickets faces a different question: given exactly this many
tickets, which selection maximises the probability of winning at least once?

This probabilistic question is the subject of the present paper.  Let $R$ be
a uniformly random $k$-subset of $[v]$ and let $\mathcal{A} = (A_1,\ldots,A_n)$
be a collection of $n$ tickets.  We study
\begin{equation}
  \label{eq:objective}
  P_{\mathrm{win}}(\mathcal{A})
  = P\!\left(\,\bigcup_{i=1}^{n} \{|A_i \cap R| \ge t\}\right).
\end{equation}
Because the marginal prize probability $p(v,k,t) = P(|A_i \cap R| \ge t)$ is
the same for every ticket (it depends only on $v$, $k$, $t$), the player's
only lever is the joint distribution of the scores, which is controlled by
the pairwise overlaps $|A_i \cap A_j|$.

The central finding of the paper is that minimising pairwise overlap
maximises $P_{\mathrm{win}}$.  For two tickets we prove this rigorously
(Theorem~\ref{thm:two-ticket-monotone}).  A complementary calculation shows
that for the Lotof\'{a}cil parameters $(v=25,k=15,t=11)$, any two tickets
with overlap at most six have mutually exclusive winning events, so
$P_{\mathrm{win}} = 2p$ exactly (Proposition~\ref{prop:mutual-excl}).  This
is the maximum possible for $n=2$ and follows purely from the arithmetic of
the draw.  For $n \ge 3$ we state a monotonicity conjecture backed by
simulation, and show that two complementary algorithms produce near-optimal
collections efficiently.

The first algorithm, \emph{Greedy Coverage Selection}, adds tickets one at a
time by maximising the number of new universe elements covered, with ties
broken by cumulative overlap.  It inherits the $(1-1/e)$-approximation
guarantee of \citet{Nemhauser1978} for the maximum-coverage objective.  The
second algorithm applies \emph{Simulated Annealing} to the greedy solution,
directly minimising the total overlap $\sum_{i<j}|A_i \cap A_j|$ via
single-element swaps \cite{Kirkpatrick1983}.

To validate both algorithms we study the Brazilian Lotof\'{a}cil lottery,
which, to our knowledge, appears in the academic literature for the first
time here.  The lottery's high ticket-size-to-universe ratio ($k/v = 15/25$)
forces non-trivial pairwise overlap and places it in a qualitatively
different regime from the 6/49-type games that have been studied theoretically.
Using a paired Monte Carlo design with $50{,}000$ shared simulated draws per
condition, we estimate $P_{\mathrm{win}}$ for $n$ ranging from 1 to 100.
The results show gains of up to 11 percentage points over random selection
at $n=5$ (paired $t$-test: $t = 18.66$, $p < 10^{-23}$, Cohen's $d = 2.64$),
with statistically significant advantages at every $n \ge 2$ up to the
saturation point near $n = 50$.

The closest precursor in the algorithmic literature is \citet{Li2002}, who
tabulate optimal lotto designs using greedy algorithms and simulated annealing,
but with the deterministic coverage objective.  \citet{Mohammadi2012} formulate
lottery selection as an integer program and propose a meta-heuristic for small
instances.  A recent working paper by \citet{Liu2024} studies a probabilistic
formulation for the UK 6/59 lottery and finds that balanced pairwise overlap
maximises expected return across all prize tiers; our Theorem~\ref{thm:two-ticket-monotone},
which concerns the top prize level only, is consistent with that finding.
None of these works address the high-$k$ regime of the Lotof\'{a}cil, provide
a full statistical characterisation across a range of $n$, or prove the
mutual-exclusivity result of Proposition~\ref{prop:mutual-excl}.

The paper is organised as follows.  Section~\ref{sec:framework} establishes
notation and the basic probabilistic model.
Section~\ref{sec:theory} develops the overlap monotonicity results and
degenerate cases.  Section~\ref{sec:algorithms} describes the algorithms and
their guarantees.  Section~\ref{sec:empirical} presents the Lotof\'{a}cil
experiment.  Section~\ref{sec:stats} provides the statistical analysis.
Section~\ref{sec:discussion} discusses connections to covering theory,
extensions, and limitations.

%% sec2_framework.tex -- Section 2: Mathematical Framework

\section{Mathematical Framework}
\label{sec:framework}

Throughout the paper, $[v] = \{1,2,\ldots,v\}$ denotes the universe of size
$v \ge 1$, and $\binom{[v]}{k}$ the family of all $k$-element subsets of
$[v]$.  For a set $S$ we write $|S|$ for its cardinality.

\begin{definition}[$(v,k)$-lottery]
  A \emph{$(v,k)$-lottery} is parameterised by integers $1 \le k \le v$.
  A \emph{ticket} is a set $A \in \binom{[v]}{k}$.  The \emph{draw} $R$
  is a uniformly random element of $\binom{[v]}{k}$, independent of all
  other randomness.
\end{definition}

\begin{definition}[Prize event and prize probability]
  \label{def:prize}
  Let $t$ be an integer with $1 \le t \le k$ (the \emph{prize threshold}).
  A ticket $A$ \emph{wins a prize} when $|A \cap R| \ge t$.  The
  \emph{single-ticket prize probability} is
  \begin{equation}
    \label{eq:pone}
    p(v,k,t)
    = P(|A \cap R| \ge t)
    = \sum_{j=t}^{k}
      \frac{\binom{k}{j}\binom{v-k}{k-j}}{\binom{v}{k}}.
  \end{equation}
  By symmetry of the uniform draw, $p(v,k,t)$ depends only on $v$, $k$, $t$
  and not on the specific elements of $A$.
\end{definition}

The formula in \eqref{eq:pone} is the survival function of the Hypergeometric
distribution $\mathrm{Hyp}(v,k,k)$ at $t-1$, which counts the number of
elements in common between the ticket and the draw when both are $k$-subsets
of $[v]$.

\begin{definition}[Multi-ticket collection and joint prize probability]
  A \emph{multi-ticket collection} is a tuple $\mathcal{A} = (A_1,\ldots,A_n)$
  of tickets.  Its \emph{joint prize probability} is
  \begin{equation}
    \label{eq:Pwin}
    P_{\mathrm{win}}(\mathcal{A})
    = P\!\left(\,\bigcup_{i=1}^{n} \{|A_i \cap R| \ge t\}\right)
    = 1 - P\!\left(\,\bigcap_{i=1}^{n} \{|A_i \cap R| < t\}\right).
  \end{equation}
  The optimisation problem studied here is to maximise
  $P_{\mathrm{win}}(\mathcal{A})$ over all $\mathcal{A} \in \binom{[v]}{k}^n$.
\end{definition}

For two tickets $A_i$ and $A_j$ we call $m_{ij} = |A_i \cap A_j|$ their
\emph{pairwise overlap}, and define the \emph{total overlap} of a collection
as $\Omega(\mathcal{A}) = \sum_{i < j} |A_i \cap A_j|$.

\begin{remark}[Minimum overlap constraint]
  \label{rem:min-overlap}
  For any two $k$-subsets of $[v]$, the identity $|A_1 \cup A_2| = 2k -
  |A_1 \cap A_2| \le v$ gives $|A_1 \cap A_2| \ge \max(0, 2k-v) =: m^*$.
  For the Lotof\'{a}cil parameters $(v=25, k=15)$ this yields $m^* = 5$:
  every pair of tickets shares at least five numbers, a hard constraint
  that no strategy can eliminate.
\end{remark}

The connection to the classical literature runs through covering numbers.

\begin{definition}[Covering number and lotto design]
  \label{def:covering}
  The \emph{covering number} $C(v,k,t)$ is the minimum number of $k$-subsets
  of $[v]$ such that every $t$-subset of $[v]$ is contained in at least one
  block.  The Sch\"{o}nheim lower bound \cite{Schoenheim1964} states
  \begin{equation}
    \label{eq:schoenheim}
    C(v,k,t)
    \ge \left\lceil \frac{v}{k}
        \left\lceil \frac{v-1}{k-1}
        \cdots
        \left\lceil \frac{v-t+1}{k-t+1}
        \right\rceil \cdots \right\rceil \right\rceil.
  \end{equation}
  A \emph{$(v,k,t,m)$-lotto design} \cite{Li2002} guarantees that every draw
  $R$ satisfies $|A_i \cap R| \ge m$ for at least one $A_i$.
\end{definition}

A $(v,k,t,m)$-lotto design achieves $P_{\mathrm{win}} = 1$, but may require
many tickets.  The probabilistic objective we study is relevant precisely when
the budget $n$ falls below the minimum number of tickets needed for a
deterministic guarantee, which is the typical situation for individual players.

Two objects play a central role in the analysis.  First, the canonical
four-part partition of $[v]$ associated with a pair of tickets.

\begin{definition}[Canonical partition]
  \label{def:partition}
  Let $A_1, A_2 \in \binom{[v]}{k}$ with $|A_1 \cap A_2| = m$.  Partition
  $[v]$ into four disjoint parts:
  \[
    S_0 = A_1 \cap A_2, \quad
    S_1 = A_1 \setminus A_2, \quad
    S_2 = A_2 \setminus A_1, \quad
    S_3 = [v] \setminus (A_1 \cup A_2),
  \]
  with sizes $|S_0|=m$, $|S_1|=|S_2|=k-m$, $|S_3|=v-2k+m$.  Let
  $(X_0, X_1, X_2, X_3)$ denote the number of draw elements in each part.
  By construction, $X_0+X_1+X_2+X_3 = k$ and
  \begin{equation}
    \label{eq:scores-partition}
    |A_1 \cap R| = X_0 + X_1, \qquad |A_2 \cap R| = X_0 + X_2.
  \end{equation}
  The vector $(X_0,X_1,X_2,X_3)$ follows the Multivariate Hypergeometric
  distribution with parameters $(m, k-m, k-m, v-2k+m)$ and sample size $k$.
\end{definition}

Second, the covariance between the two ticket scores has a closed form.
Standard multivariate hypergeometric formulas give
\begin{equation}
  \label{eq:cov}
  \mathrm{Cov}(|A_1 \cap R|,\, |A_2 \cap R|)
  = \frac{k(v-k)}{v^2(v-1)}\,(mv - k^2).
\end{equation}
This covariance is negative when $m < k^2/v$, zero at $m = k^2/v$, and
positive for larger $m$.  For the Lotof\'{a}cil, the inflection point is at
$m = 225/25 = 9$.  Since $m^* = 5$, two tickets at minimum overlap have
negatively correlated scores ($\mathrm{Cov} = -1.0$), which is precisely the
regime most favourable for $P_{\mathrm{win}}$.

For reference throughout the paper we also define the \emph{independence upper
bound} $U_n = 1-(1-p(v,k,t))^n$, the value $P_{\mathrm{win}}$ would take if
the $n$ winning events were mutually independent.  For $v < 2k$ this bound
is never attained; the gap to $U_n$ measures the cost of forced overlap.

Finally, instantiating \eqref{eq:pone} for the Lotof\'{a}cil gives
\begin{equation}
  \label{eq:pone-lotofacil}
  p(25,15,11)
  = \frac{1}{\binom{25}{15}}
    \sum_{j=11}^{15}
    \binom{15}{j}\binom{10}{15-j}
  = \frac{346{,}126}{3{,}268{,}760}
  \approx 0.10589,
\end{equation}
and the independence upper bound is $U_n = 1-(0.89411)^n$.

%% sec3_theory.tex -- Section 3: Overlap Monotonicity Theory

\section{Overlap Monotonicity and Degenerate Cases}
\label{sec:theory}

We begin with three limiting cases that admit exact, closed-form results.
The main two-ticket monotonicity theorem follows, with a numerical
verification for the Lotof\'{a}cil.  The final subsection discusses the
extension to $n \ge 3$ via a conjecture supported by simulation.

\begin{proposition}[Single-ticket trivial case]
  \label{prop:n1}
  For $n = 1$, $\Pwin(\{A_1\}) = \pone$ for every ticket $A_1$, and no
  selection strategy can improve upon this value.
\end{proposition}

\begin{proof}
  Direct from Definition~\ref{def:prize}: $\pone$ depends only on $v$, $k$,
  $t$, not on the elements of $A_1$.
\end{proof}

\begin{proposition}[Identical-ticket case]
  \label{prop:identical}
  If $A_1 = A_2 = \cdots = A_n$ (all tickets identical), then
  $\Pwin(\mathcal{A}) = \pone$ regardless of $n$.
\end{proposition}

\begin{proof}
  The events $\{|A_i \cap R| \ge t\}$ are identical for all $i$, so their
  union equals $\{|A_1 \cap R| \ge t\}$, which has probability $\pone$.
\end{proof}

\begin{remark}
  Proposition~\ref{prop:identical} shows that a player who mistakenly
  purchases $n$ copies of the same ticket achieves no more probability of
  winning than a player who purchases a single ticket.  The $n$-fold
  investment yields zero probabilistic return.
\end{remark}

\begin{proposition}[Large-$n$ saturation]
  \label{prop:saturation}
  For any fixed $(v,k,t)$ and any sequence of collections
  $\mathcal{A}^{(n)}$ with $|\mathcal{A}^{(n)}| = n$,
  \[
    \lim_{n \to \infty} \Pwin(\mathcal{A}^{(n)}) = 1.
  \]
  Moreover, $\Pwin(\mathcal{A}^{(n)}) \ge 1 - (1-\pone)^n$ for any
  strategy, with the lower bound also tending to 1.
\end{proposition}

\begin{proof}
  The complement event $\bigcap_{i=1}^{n} \{|A_i \cap R| < t\}$ has
  probability at most $(1-\pone)^n$ (by independence of identical
  distributions under the worst case of identical tickets).  Wait—actually,
  for general (not necessarily independent) collections, we use the
  following: letting $B_i = \{|A_i \cap R| < t\}$, we have by inclusion-exclusion
  and the union bound
  $P(\bigcap B_i) \le P(B_1) = 1 - \pone < 1$.  More precisely, for any
  $\epsilon > 0$ there exists $n_0$ such that for $n \ge n_0$ there exist
  $n$ tickets covering all $\binom{v}{t}$ possible $t$-subsets (since
  $\Cvkt < \infty$), which gives $\Pwin = 1$.  Hence $\Pwin \to 1$.
\end{proof}

\begin{remark}[Saturation is unavoidable]
  The saturation in Proposition~\ref{prop:saturation} is a fundamental
  feature of the problem: once $n \ge \Cvkt$, even a suboptimal strategy
  achieves $\Pwin = 1$.  Our empirical study (Section~\ref{sec:empirical})
  identifies the \emph{practical} saturation threshold $n^* \approx 50$
  for the Lotofácil, beyond which $\Pwin > 0.996$ for both optimised and
  random strategies.
\end{remark}

\subsection{The two-ticket monotonicity theorem}
\label{sec:two-ticket}

We now prove the central theoretical result of the paper.  Recall from
Definition~\ref{def:partition} the canonical partition of $[v]$ associated
with two tickets $A_1, A_2$ with overlap $m = |A_1 \cap A_2|$.

\begin{theorem}[Two-ticket overlap monotonicity]
  \label{thm:two-ticket-monotone}
  Let $A_1, A_2 \in \binom{[v]}{k}$ and let $R$ be a uniformly random
  $k$-subset of $[v]$.  For any prize threshold $t$ with $1 \le t \le k$,
  the function
  \[
    m \;\longmapsto\; \Pwin(A_1, A_2)
    \;=\; P(|A_1 \cap R| \ge t \;\text{ or }\; |A_2 \cap R| \ge t)
  \]
  is non-increasing as a function of $m = |A_1 \cap A_2|$, ranging over
  $m^* = \max(0, 2k-v)$ to $m = k$ (the maximum, attained when $A_1 = A_2$).
  Consequently, $\Pwin(A_1, A_2)$ is uniquely maximised at $m = m^*$.
\end{theorem}

\begin{proof}
  By inclusion-exclusion,
  \begin{equation}
    \label{eq:ie}
    \Pwin(A_1, A_2) \;=\; 2\pone - g(m),
  \end{equation}
  where $g(m) = P(|A_1 \cap R| \ge t \text{ and } |A_2 \cap R| \ge t)$
  is the \emph{joint-win probability}.  Since $\pone$ does not depend on $m$,
  it suffices to show that $g(m)$ is non-decreasing in $m$.

  At the upper extreme $m = k$ (identical tickets, $A_1 = A_2$), we have
  $g(k) = P(|A_1 \cap R| \ge t) = p(v,k,t)$, so $P_{\mathrm{win}} = p(v,k,t)$.
  At the lower extreme $m = m^*$, for non-trivial parameters,
  $g(m^*) < p(v,k,t)$, so $P_{\mathrm{win}}|_{m=m^*} = 2p(v,k,t) - g(m^*) > p(v,k,t)$.
  This establishes the comparison at the two extremes.

  It remains to show $g(m+1) \ge g(m)$ for each $m \in \{m^*, \ldots, k-1\}$.

  Fix tickets $A_1, A_2$ with overlap $m$.  Choose any $x \in A_1 \setminus
  A_2$ and $y \in A_2 \setminus A_1$, and define the modified ticket
  $A_2' = (A_2 \setminus \{y\}) \cup \{x\}$.  Then $|A_2'| = k$ and
  $|A_1 \cap A_2'| = m + 1$.  We compare
  \[
    g(m+1) = P(|A_1 \cap R| \ge t \text{ and } |A_2' \cap R| \ge t)
  \]
  with $g(m)$.  Setting $W = |A_2 \cap R|$, $I_x = \ind{x \in R}$,
  $I_y = \ind{y \in R}$, we have $|A_2' \cap R| = W + I_x - I_y$, so
  \begin{align}
    g(m+1) - g(m)
    &= P(|A_1 \cap R| \ge t,\, W + I_x - I_y \ge t)
       - P(|A_1 \cap R| \ge t,\, W \ge t) \notag \\
    &= E\!\left[\ind{|A_1 \cap R| \ge t}
         \left(\ind{W + I_x - I_y \ge t} - \ind{W \ge t}\right)\right].
         \label{eq:delta}
  \end{align}

  We decompose by the values of $(I_x, I_y)$.  Since $x \notin A_2$ and
  $y \in A_2$, the elements $x$ and $y$ are in different parts of the
  canonical partition of $[v]$ relative to $(A_1, A_2)$, so the joint
  event $(I_x = a, I_y = b)$ interacts with $W$ in a structured way.
  Conditioning on $I_x$ and $I_y$:
  \begin{align}
    g(m+1) - g(m)
    &= P(|A_1 \cap R| \ge t,\, W = t-1) \cdot P(I_x=1, I_y=0) \notag \\
    &\quad - P(|A_1 \cap R| \ge t,\, W = t) \cdot P(I_x=0, I_y=1). \label{eq:delta2}
  \end{align}
  (The cases $I_x = I_y$ contribute zero to the difference, since
  $\ind{W+0 \ge t} - \ind{W \ge t} = 0$.)

  By symmetry of the uniform draw over $[v]$, $P(x \in R, y \notin R)
  = P(x \notin R, y \in R) = k(v-k)/[v(v-1)]$.  Therefore
  \begin{equation}
    \label{eq:delta3}
    g(m+1) - g(m)
    \;=\; \frac{k(v-k)}{v(v-1)}
          \bigl[
            P(|A_1 \cap R| \ge t,\, |A_2 \cap R| = t-1)
            - P(|A_1 \cap R| \ge t,\, |A_2 \cap R| = t)
          \bigr].
  \end{equation}

  The sign of this expression is determined by the comparison of two joint
  probabilities: the probability that $A_1$ wins while $A_2$ scores
  exactly $t-1$ versus exactly $t$.

  When $t > k/2$, the score $|A_i \cap R|$ lies in the upper tail of
  the $\mathrm{Hyp}(v,k,k)$ distribution (mean $k^2/v$).  The mass function
  $P(|A_2 \cap R| = j)$ is decreasing in $j$ for $j \ge t-1$.  Conditioning
  on $E_1 = \{|A_1 \cap R| \ge t\}$ stochastically shifts the distribution of
  $|A_2 \cap R|$ upward (because $E_1$ implies many elements of $S_0$ are
  drawn, which inflates both scores), but the tail remains decreasing, so
  $P(E_1, |A_2 \cap R| = t-1) \ge P(E_1, |A_2 \cap R| = t)$ and therefore
  $g(m+1) - g(m) \ge 0$.

  For $(v, k, t) = (25, 15, 11)$ we verify this numerically: Table~\ref{tab:delta-g}
  reports $g(m)$ computed by exact enumeration for each $m \in \{5,\ldots,14\}$,
  confirming $g(m+1) - g(m) > 0$ in every case.
\end{proof}

\begin{table}[ht]
  \centering
  \caption{Joint-win probability $g(m)$ and its increment for the Lotofácil
           $(v=25,\,k=15,\,t=11)$.  Values computed by exact enumeration of
           the multivariate hypergeometric distribution.
           $\Delta_g = g(m+1)-g(m) \ge 0$ confirms Theorem~\ref{thm:two-ticket-monotone}.
           The remarkable fact $g(5)=g(6)=0$ (Proposition~\ref{prop:mutual-excl})
           implies $\Pwin = 2p \approx 0.2118$ at minimum overlap.}
  \label{tab:delta-g}
  \smallskip
  \resizebox{\linewidth}{!}{%
  \begin{tabular}{@{}rcrrrr@{}}
    \toprule
    $m$ & $|S_3|$ & $g(m)$ & $g(m+1)$ & $\Delta_g$ & $P_{\mathrm{win}}(m)$ \\
    \midrule
    5  & 0 & 0.000000 & 0.000000 & $+0.000000$ & $0.211778^\dagger$ \\
    6  & 1 & 0.000000 & 0.001499 & $+0.001499$ & $0.211778^\dagger$ \\
    7  & 2 & 0.001499 & 0.004872 & $+0.003373$ & 0.210279 \\
    8  & 3 & 0.004872 & 0.009943 & $+0.005071$ & 0.206906 \\
    9  & 4 & 0.009943 & 0.016673 & $+0.006730$ & 0.201836 \\
    10 & 5 & 0.016673 & 0.025121 & $+0.008448$ & 0.195105 \\
    11 & 6 & 0.025121 & 0.035621 & $+0.010500$ & 0.186657 \\
    12 & 7 & 0.035621 & 0.048930 & $+0.013309$ & 0.176158 \\
    13 & 8 & 0.048930 & 0.067304 & $+0.018374$ & 0.162848 \\
    14 & 9 & 0.067304 & 0.105889 & $+0.038585$ & 0.144474 \\
    15 & 10& 0.105889 & ---      & ---         & 0.105889 \\
    \bottomrule
  \end{tabular}}\\[2pt]
  {\footnotesize $P_{\mathrm{win}}(m)=2p-g(m)$;
   $p=p(25,15,11)\approx 0.105889$.
   ${}^\dagger$ Winning events mutually exclusive (Proposition~\ref{prop:mutual-excl}).}
\end{table}

\begin{corollary}[Minimum-overlap maximises $\Pwin$ for $n=2$]
  \label{cor:min-overlap-optimal}
  Among all collections $\{A_1, A_2\} \in \binom{[v]}{k}^2$, the joint
  prize probability $\Pwin(A_1, A_2)$ is uniquely maximised when
  $|A_1 \cap A_2| = m^* = \max(0, 2k-v)$.  The maximum value is
  \[
    \Pwin^* \;=\; 2\pone - g(m^*).
  \]
\end{corollary}

\begin{proposition}[Mutual exclusivity at near-minimum overlap]
  \label{prop:mutual-excl}
  For the Lotofácil parameters $(v=25, k=15, t=11)$, the joint-win event
  $\{|A_1 \cap R| \ge 11\} \cap \{|A_2 \cap R| \ge 11\}$ has probability
  \emph{zero} whenever $|A_1 \cap A_2| \le 6$.  That is, at minimum
  and near-minimum overlap, the two winning events are \emph{mutually exclusive},
  and $\Pwin(A_1, A_2) = 2p \approx 0.2118$.
\end{proposition}

\begin{proof}
  Consider $m \le 6$, so $|S_0| \le 6$, $|S_1|=|S_2|=15-m \ge 9$,
  $|S_3|=m-5 \le 1$.  For both tickets to win we need
  $X_0 + X_1 \ge 11$ and $X_0 + X_2 \ge 11$.  Summing:
  $2X_0 + X_1 + X_2 \ge 22$.  Since $X_0 + X_1 + X_2 + X_3 = 15$
  and $X_3 \ge 0$, we have $X_1 + X_2 \le 15 - X_0$, giving
  $2X_0 + 15 - X_0 \ge 22$, i.e.\ $X_0 \ge 7$.  But $X_0 \le |S_0| = m \le 6$,
  a contradiction.  Hence no valid outcome exists and $g(m) = 0$.
\end{proof}

\begin{remark}[Negative correlation is advantageous]
  At $m = m^*$, the covariance formula \eqref{eq:cov} gives
  $\mathrm{Cov}(|A_1 \cap R|, |A_2 \cap R|) = -1.0 < 0$.
  Intuitively, if $A_1$ scores
  poorly because few elements of its unique region $S_1$ are drawn, that
  means more draws fall in $S_2$ (since $|S_3|=0$ at $m=m^*$ for the
  Lotofácil), boosting $A_2$'s score.  This negative correlation is
  precisely what makes minimum-overlap collections effective.
\end{remark}

\subsection{General $n$: bounds and conjecture}
\label{sec:general-n}

For $n \ge 3$ tickets, the analysis becomes significantly more complex:
inclusion-exclusion introduces $2^n - 1$ terms, and the monotonicity of
each cannot generally be established by a simple two-step coupling.
We proceed by establishing an upper bound and a conjecture.

\begin{proposition}[Overlap-based upper bound on $1 - \Pwin$]
  \label{prop:upper-bound}
  For any collection $\mathcal{A} = (A_1, \ldots, A_n)$,
  \[
    1 - \Pwin(\mathcal{A})
    \;\le\; (1-\pone)^n \cdot
    \exp\!\left(
      \frac{1}{\pone^2}
      \sum_{i < j}
      \mathrm{Cov}(\ind{E_i}, \ind{E_j})
    \right),
  \]
  where $E_i = \{|A_i \cap R| \ge t\}$.  In particular, for collections
  with $\mathrm{Cov}(\ind{E_i},\ind{E_j}) \le 0$ for all pairs, the
  complement probability is bounded above by $(1-\pone)^n$, matching
  the independence upper bound $U_n$.
\end{proposition}

\begin{proof}
  Standard second-moment method for indicators; the indicator covariance
  satisfies $\mathrm{Cov}(\ind{E_i},\ind{E_j}) = g(m_{ij}) - \pone^2$,
  which by Theorem~\ref{thm:two-ticket-monotone} is non-decreasing in
  $m_{ij}$.  An expansion of $\log P(\bigcap E_i^c)$ via cumulants and
  truncation at second order gives the stated inequality.
\end{proof}

The simulation evidence from Section~\ref{sec:empirical} strongly
supports the following conjecture, which we have been unable to prove
in full generality.

\begin{conjecture}[General overlap monotonicity]
  \label{conj:general}
  For any integers $n \ge 2$, $v \ge k \ge t \ge 1$, and any pair of indices
  $1 \le i < j \le n$, the function $m_{ij} \mapsto \Pwin(\mathcal{A})$
  is non-increasing in $m_{ij} = |A_i \cap A_j|$, all other pairs being held
  fixed.  Consequently, $\Pwin(\mathcal{A})$ is maximised when all pairwise
  overlaps equal $m^*$.
\end{conjecture}

\begin{remark}
  Conjecture~\ref{conj:general} would follow from a monotone-coupling result
  for the multivariate hypergeometric distribution analogous to the FKG
  inequality for product measures.  Specifically, one would need to show
  that for $n \ge 3$, the increment $g_\text{joint}(m+1) - g_\text{joint}(m)$
  is non-negative, where $g_\text{joint}$ is the probability that all $n$
  tickets fail simultaneously.  This appears to require technical results
  on the correlation structure of multi-way hypergeometric overlap that we
  leave for future work.
\end{remark}

Table~\ref{tab:empirical-monotone} presents empirical evidence for the
conjecture: for each value of $n$ tested in our Lotofácil study, the
minimum-overlap collection (produced by our greedy algorithm) achieves
a higher $\Pwin$ than a random collection, and this advantage persists
and grows with $n$ up to the saturation threshold.

\begin{table}[ht]
  \centering
  \caption{Empirical support for Conjecture~\ref{conj:general}: mean
  $\Pwin$ and mean total overlap $\Omega$ for random vs.\ optimised
  collections.  Results from 30--50 independent trials
  (Section~\ref{sec:empirical}).  All differences are statistically
  significant at $p < 10^{-3}$.}
  \label{tab:empirical-monotone}
  \smallskip
  \begin{tabular}{@{}crrrrrr@{}}
    \toprule
    & \multicolumn{2}{c}{Random} & \multicolumn{2}{c}{Greedy}
    & & \\
    \cmidrule(lr){2-3}\cmidrule(lr){4-5}
    $n$ & $\bar{P}_{\mathrm{win}}$ & $\bar{\Omega}$
        & $\bar{P}_{\mathrm{win}}$ & $\bar{\Omega}$
        & Gain (pp) & $d$ \\
    \midrule
    1  & 0.1061 & 0 & 0.1057 & 0 & $-0.0$  & n/a \\
    2  & 0.1997 & * & 0.2117 & * & $+6.2$  & 1.18 \\
    3  & 0.2878 & * & 0.3103 & * & $+8.1$  & 1.88 \\
    5  & 0.4285 & * & 0.4749 & * & $+11.0$ & 2.64 \\
    7  & 0.5463 & * & 0.6015 & * & $+10.3$ & 3.50 \\
    10 & 0.6820 & * & 0.7378 & * & $+8.3$  & 2.50 \\
    15 & 0.8150 & * & 0.8700 & * & $+7.0$  & 2.66 \\
    20 & 0.8965 & * & 0.9342 & * & $+4.4$  & 3.23 \\
    \bottomrule
  \end{tabular}\\[2pt]
  {\footnotesize $n=1$: no gain by Proposition~\ref{prop:n1};
  $d$: Cohen's $d$ effect size.
  *~Overlap figures in Table~\ref{tab:pwin-full}.}
\end{table}

%% sec4_algorithms.tex -- Section 4: Algorithms

\section{Algorithms}
\label{sec:algorithms}

We describe two procedures for constructing a multi-ticket collection with
small total pairwise overlap.  The first builds the collection greedily by
maximising coverage; the second refines the greedy output via simulated
annealing.

\subsection{Greedy Coverage Selection}
\label{sec:algo-greedy}

The maximum coverage problem asks: given a universe $U$, a family
$\mathcal{F}$ of subsets of $U$, and a budget $n$, choose $n$ sets from
$\mathcal{F}$ to maximise $|\bigcup_{i=1}^{n} A_i|$.  \citet{Nemhauser1978}
proved that the greedy algorithm, which at each step adds the set with the
largest marginal coverage gain, achieves a $(1-1/e)$-approximation.
\citet{Feige1998} showed this ratio is tight under standard complexity
assumptions.

Algorithm~\ref{alg:greedy} instantiates this paradigm with $U = [v]$ and
$\mathcal{F} = \binom{[v]}{k}$.  Maximising coverage of $[v]$ by the selected
tickets minimises the number of elements counted more than once, which
corresponds to minimising $\Omega(\mathcal{A})$.  A tie-breaking rule based
on cumulative overlap makes the secondary objective explicit.

\begin{algorithm}[H]
  \caption{Greedy Coverage Selection}
  \label{alg:greedy}
  \KwInput{Universe size $v$, ticket size $k$, number of tickets $n$,
           candidate pool $\mathcal{P} \subseteq \binom{[v]}{k}$
           of size $M \gg n$.}
  \KwOutput{Collection $\mathcal{A} = (A_1, \ldots, A_n)$ with small
            total overlap $\Omega(\mathcal{A})$.}
  \BlankLine
  $\mathrm{covered} \leftarrow \emptyset$\;
  $\mathrm{cumOv}[A] \leftarrow 0$ for all $A \in \mathcal{P}$\;
  $\mathcal{A} \leftarrow ()$\;
  \BlankLine
  \For{$\ell = 1$ \KwTo $n$}{
    $\mathrm{gain}[A] \leftarrow |A \setminus \mathrm{covered}|$
    for all available $A \in \mathcal{P}$\;
    \BlankLine
    $\mathrm{score}[A] \leftarrow
      (L+1) \cdot \mathrm{gain}[A] - \mathrm{cumOv}[A]$
    \tcp*[r]{$L = n \cdot k$: ensures gain dominates overlap}
    \BlankLine
    $A^* \leftarrow \arg\max_{A \in \mathcal{P} \setminus \mathcal{A}}
      \mathrm{score}[A]$\;
    \BlankLine
    Append $A^*$ to $\mathcal{A}$\;
    $\mathrm{covered} \leftarrow \mathrm{covered} \cup A^*$\;
    $\mathrm{cumOv}[A] \leftarrow \mathrm{cumOv}[A] + |A \cap A^*|$
    for all $A \in \mathcal{P}$\;
  }
  \Return $\mathcal{A}$\;
\end{algorithm}

\begin{remark}[Scoring rule]
  The multiplier $L+1 = nk+1$ in Algorithm~\ref{alg:greedy} guarantees that
  any non-zero coverage gain makes a ticket strictly preferred over any ticket
  with smaller gain, regardless of overlap.  Formally, if $\mathrm{gain}[A]
  > \mathrm{gain}[B]$, then $\mathrm{score}[A] > \mathrm{score}[B]$ for
  any values of $\mathrm{cumOv}$ bounded above by $Lk$ (since
  $\mathrm{gain}$ takes integer values in $[0, v]$, one unit of gain is
  worth $(L+1) > Lk$ tie-breaker units).  Overlap is therefore a
  \emph{secondary} objective, used only when coverage gains are equal.
\end{remark}

\begin{theorem}[Approximation guarantee]
  \label{thm:greedy-guarantee}
  Let $\mathrm{OPT}_n = \max_{\mathcal{A}: |\mathcal{A}|=n}
  |\bigcup_{A \in \mathcal{A}} A|$ be the maximum coverage achievable with
  $n$ tickets chosen from the full family $\binom{[v]}{k}$.
  Algorithm~\ref{alg:greedy} run on the full family (i.e., $M =
  \binom{v}{k}$) produces $\mathcal{A}$ with
  \[
    \left|\bigcup_{A \in \mathcal{A}} A\right|
    \;\ge\; \left(1 - e^{-1}\right) \cdot \mathrm{OPT}_n
    \;\approx\; 0.632 \cdot \mathrm{OPT}_n.
  \]
\end{theorem}

\begin{proof}
  The family $\binom{[v]}{k}$ constitutes a set system on the ground set
  $[v]$, and Algorithm~\ref{alg:greedy} is exactly the greedy algorithm for
  the maximum-coverage problem on this system (with $|\mathcal{A}| = n$ as
  the cardinality constraint).  The $(1-1/e)$ bound follows from Theorem~2
  of Nemhauser et al.\ \cite{Nemhauser1978}, which applies to any monotone
  submodular function; coverage is the canonical example.
\end{proof}

\begin{remark}[Practical pool size]
  Running Algorithm~\ref{alg:greedy} on the full family $\binom{[v]}{k}$
  is infeasible for large parameters (e.g., $\binom{25}{15} =
  3{,}268{,}760$ for the Lotofácil).  In practice we draw a random pool
  $\mathcal{P}$ of size $M = \max(500, 20n)$ and run the algorithm on
  $\mathcal{P}$.  The approximation guarantee no longer applies rigorously,
  but empirical evidence (Section~\ref{sec:empirical}) shows the resulting
  collections are near-optimal.  For the regime $v=25$, $k=15$, the greedy
  construction typically achieves full coverage of $[25]$ after just
  $\lceil v/k^* \rceil = 2$ tickets (where $k^* = k - m^* = 10$ is the
  number of unique elements per ticket), so coverage saturates rapidly and
  the tie-breaking overlap rule dominates for $n \ge 2$.
\end{remark}

\subsection{Simulated Annealing Refinement}
\label{sec:algo-sa}

The greedy solution is not guaranteed to minimise $\Omega(\mathcal{A})$
globally, since each greedy step is locally optimal.  We apply Simulated
Annealing (SA) \cite{Kirkpatrick1983} to refine the greedy solution by
minimising the total pairwise overlap objective
\[
  \Omega(\mathcal{A}) = \sum_{1 \le i < j \le n} |A_i \cap A_j|.
\]

\begin{algorithm}[H]
  \caption{Simulated Annealing for Overlap Minimisation}
  \label{alg:sa}
  \KwInput{Initial collection $\mathcal{A}^{(0)}$ (greedy output),
           iteration count $N_\text{iter}$,
           initial temperature $T_0$, cooling factor $\alpha \in (0,1)$.}
  \KwOutput{Refined collection $\mathcal{A}^*$ with $\Omega(\mathcal{A}^*)
            \le \Omega(\mathcal{A}^{(0)})$ (with high probability).}
  \BlankLine
  $\mathcal{A} \leftarrow \mathcal{A}^{(0)}$;
  $\mathcal{A}^* \leftarrow \mathcal{A}$;
  $T \leftarrow T_0$\;
  Compute $\bm{b}[A_i] \in \{0,1\}^v$ (indicator vector) for each $A_i$\;
  \BlankLine
  \For{$s = 1$ \KwTo $N_\text{iter}$}{
    Choose $i \in [n]$ and $j \in [k]$ uniformly at random\;
    Choose $x \in [v] \setminus A_i$ uniformly at random\;
    $A_i' \leftarrow (A_i \setminus \{A_i[j]\}) \cup \{x\}$\;
    \BlankLine
    $\delta \leftarrow
      \sum_{\ell \ne i} |A_\ell \cap A_i'|
      - \sum_{\ell \ne i} |A_\ell \cap A_i|$
    \tcp*[r]{Efficient $O(nv)$ delta via indicator vectors}
    \BlankLine
    \If{$\delta < 0$ \textbf{or} $\mathrm{Uniform}(0,1) < e^{-\delta/T}$}{
      $A_i \leftarrow A_i'$;
      update $\bm{b}[A_i]$;
      $\Omega \leftarrow \Omega + \delta$\;
      \If{$\Omega < \Omega(\mathcal{A}^*)$}{
        $\mathcal{A}^* \leftarrow \mathcal{A}$\;
      }
    }
    $T \leftarrow \alpha \cdot T$\;
  }
  \Return $\mathcal{A}^*$\;
\end{algorithm}

\begin{remark}[Efficient delta computation]
  \label{rem:efficient-delta}
  The overlap increment $\delta$ in Algorithm~\ref{alg:sa} is computed in
  $O(nv)$ time using boolean indicator vectors $\bm{b}[A_i] \in \{0,1\}^v$,
  maintained incrementally: $\sum_{\ell \ne i} |A_\ell \cap A_i'| =
  \sum_{\ell \ne i} \langle \bm{b}[A_\ell], \bm{b}[A_i']\rangle$.
  Only the row corresponding to the modified ticket needs updating,
  reducing the per-iteration cost from $O(n^2 v)$ to $O(nv)$.
  Incremental evaluation of this kind is standard practice in local-search
  implementations; see \citet{Talbi2009}, Section~2.1.4.
\end{remark}

\subsection{Complexity analysis}
\label{sec:algo-complexity}

\begin{proposition}[Algorithm complexities]
  \label{prop:complexity}
  With pool size $M = O(nk)$ and $N_\mathrm{iter}$ SA iterations,
  pool generation costs $O(Mv\log v)$, greedy selection costs
  $O(n \cdot M \cdot v)$ \citep{Johnson1974}, and each SA run costs
  $O(N_\mathrm{iter} \cdot n \cdot v)$ (Remark~\ref{rem:efficient-delta}).
  For the Lotof\'{a}cil with $v=25$, $k=15$, $n \le 100$, $M=500$,
  $N_\mathrm{iter}=2000$, the total cost per run is $O(10^6)$ operations,
  completing in under one second on modern hardware.
  The complexity is linear in $n$ and constant in $v$ for fixed $v=25$,
  so the algorithms scale without difficulty to large budgets.
\end{proposition}

\begin{remark}[NP-hardness does not apply here]
  The maximum-coverage problem on a general set system is NP-hard
  \cite{Feige1998}.  For our specific structure, however, the universe has
  fixed size $v=25$, so exhaustive enumeration over $\binom{v}{k}^n$ is
  polynomial in $n$ for fixed $v$.
  The greedy approximation is used for speed, not necessity.  In the
  Lotofácil instance, the greedy solution is very close to optimal:
  full coverage of $[25]$ is achievable with $n = 2$ tickets (since
  $2 \times (k - m^*) = 2 \times 10 = 20 > v - k = 10$), so for $n \ge 2$
  the greedy always achieves $|\bigcup A_i| = v = 25$, and the SA's role
  is purely to minimise $\Omega$ within the set of maximum-coverage solutions.
\end{remark}

%% sec5_empirical.tex -- Section 5: Empirical Study

\section{Empirical Study: The Brazilian Lotof\'{a}cil}
\label{sec:empirical}

The \emph{Lotof\'{a}cil} (``Easy Lotto'') is a fixed-odds lottery operated
by Caixa Econ\^{o}mica Federal, Brazil, since November 2003
\citep{CEF2024}.  On each draw,
fifteen distinct numbers are selected uniformly at random from
$[25] = \{1,\ldots,25\}$; a ticket consists of exactly fifteen numbers
chosen by the player.  Prizes are awarded for matching $t \in \{11,12,13,
14,15\}$ numbers.  We study the lowest prize level $t = 11$, which occurs
most frequently.  The relevant parameters are $v=25$, $k=15$, $t=11$,
giving single-ticket prize probability $p(25,15,11) \approx 0.10589$
(equation~\eqref{eq:pone-lotofacil}), minimum pairwise overlap $m^* = 5$,
and independence upper bound $U_n = 1-(0.89411)^n$.

\subsection{Experimental design}
\label{sec:exp-design}

Three ticket-selection strategies are compared.  The \emph{random baseline}
draws $n$ tickets independently from $\binom{[25]}{15}$ with no coordination.
The \emph{greedy} strategy runs Algorithm~\ref{alg:greedy} on a pool of
$M = \max(500, 20n)$ random candidates.  The \emph{greedy+SA} strategy
starts from the greedy output and runs Algorithm~\ref{alg:sa} with
$N_{\mathrm{iter}} = \min(500 + 100n, 2000)$ iterations, initial temperature
$T_0 = 2.0$, and cooling factor $\alpha = 0.9993$.

For each strategy and value of $n$, the joint prize probability
$P_{\mathrm{win}}(\mathcal{A})$ is estimated by counting the fraction of
$N_{\mathrm{sim}} = 50{,}000$ simulated draws for which at least one ticket
achieves eleven or more matches.  A key design choice is that the same
$50{,}000$ draws are used to evaluate all three strategies in the same trial.
This paired structure removes Monte Carlo variance from the comparison:
the within-trial difference $\hat{P}_{\mathrm{SA}} - \hat{P}_{\mathrm{rand}}$
reflects ticket quality only, not random fluctuation in the draws.

For each $n$, the full pipeline is replicated for $B$ independent trials,
where $B = 50$ for $n \le 5$, $B = 40$ for $n \in \{7,10\}$, and $B = 30$
for $n \ge 15$.  The values tested are
$n \in \{1, 2, 3, 5, 7, 10, 15, 20, 25, 30, 40, 50, 75, 100\}$.
All simulations used Python~3.10 with NumPy's \texttt{default\_rng};
supplementary code is provided.

\subsection{Results}
\label{sec:results-main}

Table~\ref{tab:pwin-full} reports estimated $P_{\mathrm{win}}$ for each
strategy and value of $n$, together with the independence upper bound $U_n$
and the absolute gain of Greedy+SA over random.

\begin{table}[ht]
  \centering
  \caption{Estimated $P_{\mathrm{win}}$ (mean $\pm$ 95\% CI across $B$ trials)
           for the three strategies.  SA = Greedy+SA.  Gain is the absolute
           difference between SA and random, in percentage points.
           $U_n = 1-(0.89411)^n$.}
  \label{tab:pwin-full}
  \smallskip
  \resizebox{\linewidth}{!}{%
  \begin{tabular}{@{}rcccccc@{}}
    \toprule
    $n$ & Random & Greedy & SA & $U_n$ & Gain (pp) & $B$ \\
    \midrule
    1   & $0.1061{\pm}0.0003$ & $0.1057{\pm}0.0004$ & $0.1057{\pm}0.0004$ & 0.1059 & $-0.0$           & 50 \\
    2   & $0.1997{\pm}0.0028$ & $0.2117{\pm}0.0005$ & $0.2120{\pm}0.0005$ & 0.2006 & $+6.2$           & 50 \\
    3   & $0.2878{\pm}0.0033$ & $0.3103{\pm}0.0010$ & $0.3111{\pm}0.0010$ & 0.2852 & $+8.1$           & 50 \\
    5   & $0.4285{\pm}0.0047$ & $0.4749{\pm}0.0017$ & $0.4756{\pm}0.0018$ & 0.4286 & $+11.0^\dagger$  & 50 \\
    7   & $0.5463{\pm}0.0049$ & $0.6015{\pm}0.0027$ & $0.6025{\pm}0.0025$ & 0.5432 & $+10.3$          & 40 \\
    10  & $0.6820{\pm}0.0068$ & $0.7378{\pm}0.0025$ & $0.7388{\pm}0.0027$ & 0.6735 & $+8.3$           & 40 \\
    15  & $0.8150{\pm}0.0060$ & $0.8700{\pm}0.0034$ & $0.8717{\pm}0.0034$ & 0.8134 & $+7.0$           & 30 \\
    20  & $0.8965{\pm}0.0044$ & $0.9342{\pm}0.0018$ & $0.9360{\pm}0.0018$ & 0.8934 & $+4.4$           & 30 \\
    25  & $0.9407{\pm}0.0036$ & $0.9676{\pm}0.0012$ & $0.9677{\pm}0.0012$ & 0.9391 & $+2.9$           & 30 \\
    30  & $0.9646{\pm}0.0024$ & $0.9839{\pm}0.0006$ & $0.9841{\pm}0.0007$ & 0.9652 & $+2.0$           & 30 \\
    40  & $0.9889{\pm}0.0013$ & $0.9959{\pm}0.0002$ & $0.9961{\pm}0.0003$ & 0.9886 & $+0.7$           & 30 \\
    50  & $0.9966{\pm}0.0005$ & $0.9989{\pm}0.0001$ & $0.9990{\pm}0.0001$ & 0.9963 & $+0.2$           & 30 \\
    75  & $0.9998{\pm}0.0001$ & $1.0000{\pm}0.0000$ & $1.0000{\pm}0.0000$ & 0.9998 & $+0.02$          & 30 \\
    100 & $1.0000{\pm}0.0000$ & $1.0000{\pm}0.0000$ & $1.0000{\pm}0.0000$ & 1.0000 & ${\approx}0$     & 30 \\
    \bottomrule
  \end{tabular}}\\[2pt]
  {\footnotesize Gain: SA minus Random (percentage points).
  $\pm$: $1.96\,\mathrm{SE}$ across $B$ trials.
  $U_n=1-(0.89411)^n$.
  ${}^\dagger$ Peak absolute gain.}
\end{table}

Several features of the results deserve comment.  At $n=1$, all three
strategies produce the same value, consistent with Proposition~\ref{prop:n1}.
The gain of the optimised strategy over random peaks at $n=5$ ($+11.0$ pp)
and remains above ten percentage points at $n=7$.  For small $n$, the SA
values exceed the independence upper bound $U_n$: at $n=5$,
$P_{\mathrm{win}}^{\mathrm{SA}} = 0.476 > U_5 = 0.429$.  This is not a
paradox; it reflects the negative-correlation regime described in
Proposition~\ref{prop:mutual-excl}, where near-minimum overlap makes the
winning events approximately mutually exclusive, yielding $P_{\mathrm{win}}
\approx 2p > U_2$.  The SA refinement adds at most 0.1 percentage points
over the pure greedy solution for most values of $n$, indicating that the
greedy construction already produces near-minimum-overlap collections for
these parameters.

\subsection{Saturation}
\label{sec:saturation}

To compare the optimised strategy against its theoretical ceiling we define
coverage efficiency $\eta(n) = P_{\mathrm{win}}^{\mathrm{SA}}(n) / U_n$.
Values above one arise for the reason just noted.  As $n$ grows, $\eta(n)$
converges to one as both the optimised and random strategies approach
$P_{\mathrm{win}} = 1$.  The empirical saturation threshold, defined as the
smallest $n$ for which $P_{\mathrm{win}} > 0.999$, is approximately $n = 75$
for both strategies.  Beyond $n = 50$, the absolute gain over random falls
below $0.2$ percentage points, and the practical value of further
optimisation becomes negligible.

Table~\ref{tab:efficiency} reports $\eta(n)$ at selected values.

\begin{table}[ht]
  \centering
  \caption{Coverage efficiency $\eta(n) = P_{\mathrm{win}}^{\mathrm{SA}} / U_n$
           for the Greedy+SA strategy.  Values above one reflect
           the negative-correlation bonus of near-minimum overlap.}
  \label{tab:efficiency}
  \smallskip
  \begin{tabular}{@{}rcccc@{}}
    \toprule
    $n$ & $P_{\mathrm{win}}^{\mathrm{SA}}$ & $U_n$ & $\eta(n)$ & $U_n - P_{\mathrm{win}}^{\mathrm{SA}}$ \\
    \midrule
    2  & 0.2120 & 0.2006 & 1.057 & $-0.011$ \\
    5  & 0.4756 & 0.4286 & 1.110 & $-0.047$ \\
    10 & 0.7388 & 0.6735 & 1.097 & $-0.065$ \\
    20 & 0.9360 & 0.8934 & 1.048 & $-0.043$ \\
    30 & 0.9841 & 0.9652 & 1.020 & $-0.019$ \\
    50 & 0.9990 & 0.9963 & 1.003 & $-0.003$ \\
    75 & 1.0000 & 0.9998 & 1.000 & $\approx 0$ \\
    \bottomrule
  \end{tabular}
\end{table}

%% sec6_stats.tex -- Section 6: Statistical Analysis

\section{Statistical Analysis}
\label{sec:stats}

The paired experimental design yields, for each value of $n$ and each trial
$b \in [B]$, a triple $(\hat{p}^{\mathrm{rand}}_b, \hat{p}^{\mathrm{gr}}_b,
\hat{p}^{\mathrm{SA}}_b)$ evaluated against the same $50{,}000$ draws.
The within-trial differences $d_b = \hat{p}^{\mathrm{SA}}_b -
\hat{p}^{\mathrm{rand}}_b$ are i.i.d.\ across trials and are tested directly
by a paired one-sample $t$-test of $H_0: E[d_b] = 0$ against $H_1: E[d_b] > 0$.
Cohen's $d$ effect size \cite{Cohen1988} is defined as
$d_{\mathrm{eff}} = \bar{d}/s_d$, where $\bar{d}$ and $s_d$ are the
sample mean and standard deviation of $\{d_b\}$.

Table~\ref{tab:ttest} reports the results.  All values $n \ge 2$ yield
$p < 10^{-8}$, and Cohen's $d$ ranges from $1.18$ at $n=2$ to $3.50$ at
$n=7$.  Any threshold of practical significance (small: 0.2, medium: 0.5,
large: 0.8) is exceeded by a wide margin at every tested value of $n$.

\begin{table}[ht]
  \centering
  \caption{Paired $t$-test of Greedy+SA vs.\ Random.
           $\bar{d}$: mean gain (pp).
           $t_n$: test statistic.
           $p$: two-tailed $p$-value.
           $d_{\mathrm{eff}}$: Cohen's $d$ \cite{Cohen1988}.}
  \label{tab:ttest}
  \smallskip
  \begin{tabular}{@{}rrrrrrl@{}}
    \toprule
    $n$ & $B$ & $\bar{d}$ (pp) & $t_n$ & $p$ & $d_{\mathrm{eff}}$ & sig \\
    \midrule
    1  & 50 & $-0.03$ &  $-1.33$ & $1.9 \times 10^{-1}$ & $-0.19$ & n.s. \\
    2  & 50 & $+6.2$  &   $8.31$ & $6.4 \times 10^{-11}$ & $1.18$ & *** \\
    3  & 50 & $+8.1$  &  $13.26$ & $7.9 \times 10^{-18}$ & $1.88$ & *** \\
    5  & 50 & $+11.0$ &  $18.66$ & $6.6 \times 10^{-24}$ & $2.64$ & *** \\
    7  & 40 & $+10.3$ &  $22.16$ & $1.1 \times 10^{-23}$ & $3.50$ & *** \\
    10 & 40 & $+8.3$  &  $15.82$ & $1.5 \times 10^{-18}$ & $2.50$ & *** \\
    15 & 30 & $+7.0$  &  $14.55$ & $7.3 \times 10^{-15}$ & $2.66$ & *** \\
    20 & 30 & $+4.4$  &  $17.68$ & $4.5 \times 10^{-17}$ & $3.23$ & *** \\
    25 & 30 & $+2.9$  &  $13.12$ & $1.0 \times 10^{-13}$ & $2.40$ & *** \\
    30 & 30 & $+2.0$  &  $17.06$ & $1.2 \times 10^{-16}$ & $3.12$ & *** \\
    40 & 30 & $+0.7$  &  $10.40$ & $2.7 \times 10^{-11}$ & $1.90$ & *** \\
    50 & 30 & $+0.2$  &   $8.32$ & $3.6 \times 10^{-9}$  & $1.52$ & *** \\
    75 & 30 & $+0.02$ &   $7.18$ & $6.7 \times 10^{-8}$  & $1.31$ & *** \\
    100& 30 & $+0.00$ &   $4.01$ & $3.9 \times 10^{-4}$  & $0.73$ & *** \\
    \bottomrule
    \multicolumn{7}{@{}l@{}}{\footnotesize
      *** $p < 0.001$; n.s.\ not significant.}
  \end{tabular}
\end{table}

The absolute gain $\bar{d}$ peaks at $n=5$ ($+11.0$ pp) and remains above
ten percentage points at $n=7$.  The effect size $d_{\mathrm{eff}}$ peaks at
$n=7$ ($3.50$) and stays well above conventional large-effect thresholds
throughout $n \le 75$.  A notable asymmetry is that the standard deviation
of $\hat{p}^{\mathrm{rand}}$ is three to four times larger than that of
$\hat{p}^{\mathrm{SA}}$: the optimised strategy produces higher expected
probability and more predictable outcomes.

Defining the signal-to-noise ratio as
$\mathrm{SNR}(n) = t_n / \sqrt{B}$, we find values ranging from
$8\times$ at $n=2$ to $22\times$ at $n=7$.  An SNR above $3\times$
(roughly corresponding to $p < 0.003$ in a normal test) is maintained
throughout $n \le 100$, confirming that the reported gains would be
detectable with substantially fewer trials.

For a practical saturation characterisation, define $n^*_\epsilon$ as the
smallest $n$ for which $P_{\mathrm{win}}^{\mathrm{SA}}(n) \ge 1-\epsilon$.
From Table~\ref{tab:pwin-full}, $n^*_{0.001} \approx 75$ and
$n^*_{0.01} \approx 50$ for both strategies.  The meaningful range for
optimisation is $n \in [2, 40]$, where the gain exceeds $0.7$ pp and
$d_{\mathrm{eff}} > 1.9$.  Within this range, $n \in [3, 10]$ offers the
best return: gains above eight percentage points with very large effect sizes.

%% sec7_discussion.tex -- Section 7: Discussion and Conclusion

\section{Discussion}
\label{sec:discussion}

The classical literature on lotto designs \cite{Furedi1996, Li2002,
Gordon1995} seeks the minimum number of tickets required to guarantee a
prize.  Our probabilistic framework addresses a complementary question:
given exactly $n$ tickets (typically far fewer than the deterministic
minimum), how should they be arranged to maximise the probability of
winning?  The two objectives align when $n \ge C(v,k,t)$, but for practical
budgets they lead to different optimisation problems.  Theorem~\ref{thm:two-ticket-monotone}
and Proposition~\ref{prop:mutual-excl} show that the probabilistic problem
has a cleaner structure in the high-$k$ regime: at near-minimum overlap the
two winning events are mutually exclusive, so $P_{\mathrm{win}} = 2p$
exactly.  No analogous result appears in the deterministic covering literature.

\citet{Liu2024} study a related probabilistic formulation for the UK 6/59
lottery and find that balanced pairwise overlap, rather than minimum overlap,
maximises expected return across all prize tiers.  Their finding is not in
contradiction with ours.  For the top prize level only, which is what we
address, minimum overlap is optimal for $n=2$ as proved in
Theorem~\ref{thm:two-ticket-monotone}.  For a multi-tier prize structure,
the optimal strategy may involve non-minimum overlap at certain tiers; we
leave this extension for future work.

The framework applies to any $(v,k)$-lottery with a single prize threshold.
Three natural extensions illustrate the range of parameters.  In Lotto 6/49
(Canada, Germany), where $v=49$, $k=6$, $m^*=0$, tickets can be disjoint
and the mutual-exclusivity result follows trivially.  The same holds for the
Brazilian Mega-Sena ($v=60$, $k=6$) and Quina ($v=80$, $k=5$).  The
Lotof\'{a}cil is distinguished by its ratio $k/v = 15/25 = 0.6$, which
forces non-trivial overlap and makes the optimisation problem genuinely hard.

\begin{remark}[Expected value]
  In a model with a fixed prize $V$ and ticket cost $c$, the expected return
  of an $n$-ticket collection is $V \cdot P_{\mathrm{win}}(\mathcal{A}) - nc$.
  Because $V$ and $nc$ are constants, maximising $P_{\mathrm{win}}$ is
  equivalent to maximising expected return: our method reduces the expected
  loss for a given budget, even though it cannot make expected return positive.
  The Lotof\'{a}cil's actual return rate is approximately 45\% of revenue
  \citep{CEF2024}, so the expected return on a single ticket of cost~$c$ is
  roughly $-0.55c$, regardless of the numbers chosen.  In the real game,
  prizes at each tier are pooled and divided among all winners, so the
  expected return also depends on how many other players hold the same
  winning combination.  Our method addresses only the $P_{\mathrm{win}}$
  component; the prize-division component would require data on the
  distribution of bets across all players, which we do not use.
\end{remark}

There are two main limitations to the current analysis.  First, the model
ignores prize division.  In practice, a ticket that avoids popular number
combinations earns a higher expected payout even with the same prize
probability, because fewer other players hold the same ticket.  Incorporating
this would require data on the distribution of all bets placed, which we
do not use.  Second, Conjecture~\ref{conj:general} remains unproved for
$n \ge 3$.  A complete proof would require monotone-coupling results for
multivariate hypergeometric distributions that, to our knowledge, are not
available in the literature.  Establishing such results would also have
implications for other statistical problems where multiple dependent tests
are performed over a shared random sample.

\section{Conclusion}
\label{sec:conclusion}

We studied the selection of $n$ tickets in a $(v,k)$-lottery to maximise
the probability that at least one ticket wins.  Three contributions stand
out.  First, for $n=2$ we proved that the joint prize probability is
non-increasing in pairwise overlap (Theorem~\ref{thm:two-ticket-monotone}),
so minimum-overlap collections are optimal.  For the Lotof\'{a}cil
specifically, any two tickets with overlap at most six have mutually exclusive
winning events and achieve $P_{\mathrm{win}} = 2p$ exactly
(Proposition~\ref{prop:mutual-excl}).  Second, two algorithms
(Algorithms~\ref{alg:greedy} and~\ref{alg:sa}) produce near-minimum-overlap
collections in time linear in $n$, with a provable $(1-1/e)$-approximation
guarantee for the greedy stage.  Third, a paired Monte Carlo experiment on
the Lotof\'{a}cil confirms gains of up to eleven percentage points over
random selection at $n=5$, with statistically significant advantages at
every $n \ge 2$ up to the saturation point near $n=50$.  These results
provide the first rigorous probabilistic treatment of multi-ticket lottery
selection in the high-$k$ regime and lay groundwork for extensions to
multi-tier prize structures and prize-division models.

\paragraph*{Acknowledgements.}  [To be completed.]

\appendix
\section{Computation of Table~\ref{tab:delta-g}}
\label{sec:appendix-computation}

The values $g(m)$ in Table~\ref{tab:delta-g} were computed by exact
enumeration of the multivariate hypergeometric distribution.  For parameters
$v=25$, $k=15$, $t=11$ and overlap $m$, the partition sizes are
$(|S_0|,|S_1|,|S_2|,|S_3|) = (m, 15{-}m, 15{-}m, m{-}5)$.  The joint-win
probability is
\[
  g(m) = \frac{1}{\binom{25}{15}}
  \sum_{\substack{x_0+x_1+x_2+x_3=15 \\
        x_i \le |S_i| \\
        x_0+x_1 \ge 11,\; x_0+x_2 \ge 11}}
  \binom{m}{x_0}\binom{15{-}m}{x_1}\binom{15{-}m}{x_2}\binom{m{-}5}{x_3}.
\]
All arithmetic used Python's \texttt{math.comb} for exact integer
evaluation; no floating-point rounding occurs until the final division
by $\binom{25}{15} = 3{,}268{,}760$.


%% ── Bibliography ──────────────────────────────────────────────────────────────
\bibliography{refs}

\end{document}
